\documentclass{article}
\usepackage[T1]{fontenc}
\usepackage{lmodern}
\usepackage{graphicx}
\usepackage{amsmath,amssymb}
\usepackage[a4paper,margin=2cm]{geometry}
\usepackage[absolute,overlay]{textpos}
\setlength{\TPHorizModule}{1mm}
\setlength{\TPVertModule}{1mm}
\usepackage[indonesian]{babel}
\usepackage{hyperref}
\setlength{\emergencystretch}{2em}
% unit: Halaman 5

\begin{document}

% === Halaman 5 (llm) ===

\begin{itemize}
    \item Salah satu metrik penting di dalam teori informasi adalah \textbf{entropi}.
    \item \textbf{Entropi} mengukur \textit{ketidakpastian} atau jumlah rata-rata informasi di dalam pesan.
    \item Semakin acak suatu data, semakin tinggi entropinya. Cipherteks adalah pesan dengan entropi yang tinggi.
    \item Entropi biasanya dinyatakan dalam satuan bit.
    \item Secara umum, entropi pesan $X$ dihitung dengan rumus:
\end{itemize}

\vspace{5mm}

\[ H(X) = - \sum_{i=1}^{n} p(x_i) \log_2 p(x_i) \]

\vspace{5mm}

$X$ = variabel acak yang menyatakan pesan \\
$x_i$ = simbol ke-$i$ di dalam pesan \\
$n$ = banyak simbol di dalam pesan \\
$p(x_i)$ = peluang kemunculan $x_i$

\end{document}
